% Source: https://artofproblemsolving.com/community/c3h2416978_1200th_post_dedicated_to_bloggy
% Stable identity: aops-topic-2416978-q03
Let $O$ be the centre and $R$ the circumradius. Each side subtends a central angle of $30^\circ$, so by the law of cosines in triangle $Oa_1a_2$,
\[10^2=2R^2-2R^2\cos30^\circ=R^2\left(2-\sqrt3\right),\qquad R^2=\dfrac{100}{2-\sqrt3}=100\left(2+\sqrt3\right).\]
Going from $a_8$ through $a_9,a_{10},a_{11},a_{12}$ to $a_1$ uses $5$ sides, and going the other way uses $7$, so the smaller piece is the hexagon $a_8a_9a_{10}a_{11}a_{12}a_1$. It does not contain $O$, and together with triangle $Oa_8a_1$ it makes up the region $Oa_8a_9\ldots a_{12}a_1$, which consists of $5$ triangles $Oa_ka_{k+1}$.

\begin{asy}
usepackage("amsmath");
size(6cm);
defaultpen(fontsize(10pt));
real R=10/(2*sin(pi/12));
pair[] V; for(int k=0;k<12;++k) V.push(R*dir(90-30*k));
fill(V[7]--V[8]--V[9]--V[10]--V[11]--V[0]--cycle,rgb(1,0.86,0.6));
draw(V[0]--V[1]--V[2]--V[3]--V[4]--V[5]--V[6]--V[7]--V[8]--V[9]--V[10]--V[11]--cycle,linewidth(0.8));
draw(V[0]--V[7],red+linewidth(0.9));
pair O=(0,0); dot(O); label("$O$",O,E);
for(int k=7;k<12;++k) draw(O--V[k],dashed);
draw(O--V[0],dashed);
for(int k=0;k<12;++k) label("$a_{"+string(k+1)+"}$",V[k],dir(90-30*k));
label("$10$",(V[3]+V[4])/2,dir(90-105));
\end{asy}

Each triangle $Oa_ka_{k+1}$ has area $\dfrac12R^2\sin30^\circ=\dfrac{R^2}{4}$, and triangle $Oa_8a_1$ has central angle $5\cdot30^\circ=150^\circ$, so its area is $\dfrac12R^2\sin150^\circ=\dfrac{R^2}{4}$ too. Hence the smaller piece has area
\[5\cdot\dfrac{R^2}{4}-\dfrac{R^2}{4}=R^2=200+100\sqrt3\approx\boxed{373.21\text{ square feet}}.\]
(This is exactly one third of the whole dodecagon, whose area is $12\cdot\dfrac{R^2}{4}=3R^2$.)
